\documentclass[]{../../../../tex/classes/styledArticle}
\usepackage{../../../../tex/packages/hebrewSupport}
\usepackage{../../../../tex/packages/mathShortcuts}
\usepackage{../../../../tex/packages/theoremsSupport}

\usepackage{pstricks-add}

\newcommand\rim{R([a, b])}
\newcommand\daru{\ol{\Sigma}}
\newcommand\Idaru{\ol{I}}
\newcommand\dard{\underline{\Sigma}}
\newcommand\Idard{\ol{I}}

%\newcommand \bs \blacksquare

\author{שחר פרץ}
\title{\textit{חדו''א 2א} $\sim$ הרצאה 2}
\date{16 באפריל 2026}
\begin{document}
	\maketitle
	\textbf{מרצה: }יעל קרשון
	
	\subsection*{פתיחת השיעור}
	נזכר בלמות הבאות מחדו''א 1א: 
	\begin{Exercise}[תרגיל פתיחה]
		\begin{enumerate}
			\item יהיו $A, B \subset \R$ קבוצות ממשיים. נגדיר $A + B = \{a + b \mid a\in A \land b \in B\}$. הוכיחו: 
			\[ \sup (A + B) = (\sup A) + (\sup B) \]
			\item יהיו $A \subset \R$ קבוצת ממשיים, וכן $\lg > 0$. נגדיר $\lg A = \{ \lg a \mid a \in A\}$. הוכיחו $\sup(\lg A) = \lg \sup A$. 
			\item יהיו $A, B \subset \R$ קבוצות חסומות. הוכיחו: 
			\[ \sup (A \cup B) = \max\{\sup A, \sup B\} \]
			\item יהיו $A, B \subset \R$ חסומות. הוכיחו: 
			\begin{enumerate}[A.]
				\item \hfil $\sup A \le \inf B \iff \forall a \in A, b \in B\co a \le b$
				\item \hfil $\sup A \ge \inf B \iff \forall \eg > 0 \exists a \in A, b \in B \co b - a < \eg$
			\end{enumerate}
			\item $A \subset B$ חסומה, אז: 
			\[ \sup A - \inf A = \sup\{a - b \mid a, b \in A\} \]
		\end{enumerate}
	\end{Exercise}
	הרבה מהדברים שנעשה בשעתיים הקרובות ידרשו את זה, וזה יתחבא למטה. 
	
	התכונית להיום: 
	\begin{itemize}
		\item השלמות מהשיעור הקודם (\shiri{1.2}, \shiri{2.1})
		\item אינטגרביליות (\shiri{2.2.1}, \shiri{2.2.2}, אולי נספיק את \shiri{2.2.3})
	\end{itemize}
	
	\section{אינטגרל רימן}
	\subsection{דוגמאות לחישוב אינטגרל רימן}
	נתבונן בפונקציה: 
	\begin{align*}
		f \co [a, b] \to \R && f(x) = \begin{cases}
			15 & x = a \\
			5 & a < x \le b
		\end{cases}
	\end{align*}
	\theo{$\int^{b}_af = 5(b-a)$}
	
	בהגדרה הקודמת שלנו לאינטגרל רימן, נדרשנו לדעת מה הערך של האינטגרל כדי להיות מסוגלים להוכיח אינטגרביליות. 
	\begin{proof}
		יהי $\eg > 0$. יהי $\dg  = \dg_\eg = ?$ (בדיעבד אחרי כן נדע להשלים את ה־$\dg$, אך כתלות ב־$\eg$ בלבד) ותהי חלוקה מתויגת $(\Pi, \{t_i\})$ של הקטע $[a, b]$. נניח שמדד העדינות $\lg(\Pi) < \dg$. יש להוכיח $\sof{S(f, \Pi, \{t_i\}) - 5(b - a)} < \eg$. נזכר שחישבנו בעבור הפונקציה לעיל בשיעור הקודם ש־: 
		\[ \sof{S(f, \Pi, \{t_i\}) - 5(b-a)} \le 10 \cdot \lg(\Pi) < 10 \delta = \cdots \]
		מכאן, שבעבור $\dg = \frac{\eg}{10}$ נקבל: 
		\[ \cdots = \eg \]
		כנדרש. 
	\end{proof}
	
	נזכר בתרגיל (''שאתם לא מבינים כמה הוא חשוב וכמה סטודנטים מתקשים בו``. ''אם אתם מתקשים, תבואו לשעת הקבלה. אם אתם לא מתקשים, תבואו לשאול משהו אחר. לא ממליצה לשאול את ה־AI, הוא יכניס לכם רעש``). יעל מזכירה ש־AI יכול מאוד לעזור, אבל צריך להיות במצב (מתמטית) שאנו יודעים להנחות אותו לעזור לנו. התרגיל – השלילה של הביטוי ''$f$ איננה אינטגרבילית``. 
	
	\theo{פונקציית דיריכלה איננה אינטגרבילית, על אף תת־קטע סגור $[a, b]$ ($a < b$). }
	\begin{proof}
		נזכר בהגדרתה: 
		\[ D(x) = \begin{cases}
			1 & x \in \Q \\
			0 & x \notin \Q
		\end{cases} \]
		נקבע קטע $[a, b]$. נניח בשלילה שקיים $I \in \R$  כך ש־$\int^b_a D(x) \dx = I$. נבחר $\eg = ?$ (נשלים אחרי כן מה הוא, כתלות ב־$I$ ו־$[a, b]$ בלבד) ותהי $\dg = \dg_e > 0$ המתאימה ל־$\eg$ ע''פ הגדרת אינטגרל רימן. נבחר $\Pi$ חלוקה של הקטע $[a, b]$, עם מדד עדינות $\lg(\Pi)$ (תרגיל: הוכיחו שקיימת כזו בהכרח כזו). מבחירת $\dg$, לכל תיוג של $\Pi$ נאמר $\{t_i\}$, מתקיים ש־: $\sof{S(f, \Pi, \{t_i\}) - I} < \eg$. מצפיפות הרציונליים והאי־רציונליים:
		\begin{itemize}
			\item קיים תיוג $\{t_i\}$ עם $t_i \in \Q$ (עבורו: סכום רימן הוא $b - a$, הוכחנו בשיעור הקודם)
			\item קיים תיוג $\{t'_i\}$ עם $t_i \notin \Q$ (עבורו: סכום רימן הוא $0$ הוכחנו בשיעור הקודם)
		\end{itemize}
		עבור התיוג הראשון נקבל $\sof{(b - a) - I} < \eg$ ועבור התיוג השני נקבל $\sof{0 - I} < \eg$. נבחר בדיעבד (וכתלות ב־$[a, b]$ בלבד) את $\eg = \frac{b - a}{3}$. מא''ש המשולש: 
		\[ \sof{b - a} \le \sof{(b - a) - I} + \sof I < 2 \eg < b - a\quad \bot \]\envendproof
	\end{proof}
	
	כאן היו אמורים לבוא דיבורים על פונקציות קדומות, שכדי לא לצבור פער (וכי הכל ברשימות ובתרגול ממילא) נדלג עליהם. 
	
	
	\subsection{חסימות ואינטגרביליות}
	אנחנו קצת בבעיה. כדי להראות שפונקציה היא אינטגרבילית, היינו צריכים ''להחזיק $I$ ביד`` ולדעת מה הערך של האינטגרל. עתה נמצא ל''מסע`` כדי למצוא תנאים הכרחיים, מספיקים, ושקולים להיות פונקציה אינטגרבילית, בלי לדעת מה הוא ה־$I$. 
	
	\noti{\hfil $R([a, b]) := \{f \co [a, b] \to \R \mid \text{אינטגרבילית} f\}$}
	לשימושנו, אינטגרבילית=אינטגרבילית רימן. כל שאר מושגי האינטגרל שנראה בהמשך לא יקראו ''אינטגרביליות``. 
	
	\theo{נניח $f \in \rim$, אז $f$ חסומה. }
	לדוגמה, הפונקציה: 
	\begin{align*}
		f \co [0, 1] \to \R && f(x) = \begin{cases}
			\frac{1}{\sqrt x} & 0 < x \le 1 \\
			0 & x = 0
		\end{cases}
	\end{align*}
	כאשר נלמד על אינטגרל לא אמיתי, נמצא משמעות לשטח מטחת לגרף, אבל $f$ איננה אינטגרבילית רימן. נפנה להוכיח את המשפט: 
	\begin{proof}[הוכחת המשפט]
		נניח $f \co [a, b] \to \R$ אינטגרבילית רימן. מכאן שקיים $I \in \R$ כך ש־$\int^b_a f = I$. יהי $\eg$, נבחר אותו להיות $\eg = 1$ (זה לא קריטי לאיזה ערך מקבעי םאותו). תהי $\dg > 0$ מתאימה ל־$\eg$ בהתאם להגדרת אינטגרל רימן. נבחר $\Pi = \{x_i\}_{i = 0}^{n}$ חלוקה של $[a, b]$ עם מדד עדינות $\lg(\Pi) < \dg$ (זכרו, יש להוכיח שקיימת כזו, תרגיל שנשאר לבית). נקבע תיוג $\{t_i\}_{i = 1}^{n}$ כלשהו של $\Pi$. 
		
		''נחלק`` את הקטע $[a, b]$ לקטעים קטנים ונדרוש שהיא תהיה חסומה על כל אחד מקטעי הביניים בנפרד. יהי $k \in \{1, \dots n\}$, וננסה להוכיח חסימות ב־$[x_{k - 1}, x_k]$. נתבונן בתיוגים $\{t'_i\}_{i = 1}^{n}$ מהצורה: 
		\[ t'_1 \dots t'_n = \underbrace{t_1 \dots t_{k - 1}}_{\ \!\text{קבועים}}, \overbrace{t}^{\mathclap{t \in [x_{k - 1}, x_k]}}, \underbrace{t_{k + 1}, \dots, t_n}_{\ \!\text{קבועים}} \]
		מהגדרת אינטגרל רימן, סכום רימן בעבור כל תיוג בסביבת $\eg$ של $I$, כלומר: 
		\[ I - \eg < \underbrace{S(f, \Pi, \{t'_i\})}_{\mathclap{=f(t)\Delta x_k + \sum_{i \neq k}f(t_i)\Delta x_i}} < I + \eg \]
		נבחין ש־$\sum_{i \neq k} f(t_i)\Delta x_i$ קבוע ביחס לבחירת $k$, ונסמנו $C_k$. לכן לכל $t \in [x_{k - 1}, x_k]$ מתקיים: 
		\begin{gather*}
			I - \eg - C_k < f(t) \Delta x_k < T + \eg - C_k \\
			\frac{I - \eg - C_k}{\Delta x_k} < f(t) < \frac{I + \eg - C_k}{\Delta x_k}
		\end{gather*}
		''לכן $f|_{[x_{k - 1}, x_k]}$ חסומה, והטוש הזה הלך לעולמו``. 
		
		סה''כ $f$ חסומה על כמות סופית של קטעים סגורים שמכסים את $[a, b]$, כלומר היא חסומה. 
	\end{proof}
	
	שימו לב! פונקציה חסומה לא חייבת להיות אינטגרבילית, לדוגמה, תנאי דרבו. חסימות הוא תנאי הכרחי, אך לא מספיק, לאינטגרביליות. 
	
	\subsection{סכומי דרבו}
	(ציורים, להשלים אחכ)
	אם רוצים להמנע מבחירת תיוג, בהינתן חלוקה, אפשר לקחת את האינפימום והסופרמום של הפונקציה על אותו קטע חלוקה. גובה המלבנים שיבחר על ידי התיוגים יהיו חסומים ע''י זה. פורמלית:
	\defi{תהי $f \co [a, b] \to \R$ חסומה. תהי $\Pi = \{x_i\}_{i = 1}^{n}$ חלוקה של $[a, b]$, כלומר $a = x_0 < x_1 < \cdots < x_n = b$. לכל $i$, נגדיר:
	\[ M_i := \sup_{\mathclap{[x_{i - 1}, x_i]}} f \quad\quad m_i := \inf_{\mathclap{[x{i -1}, x_i]}} f \]
	(שימו לב שזה מוגדר רק בגלל שהנחנו $f$ חסומה). 
	}
	
	\fig{\begin{pspicture}(-0.5,-3)(10,3)
			\psaxes[ticks=none,labels=none]{->}(0,0)(-1,-2.5)(10,3)
			\psplot[plotpoints=100,linewidth=1.5pt,algebraic]{0}{9.5}{sqrt(x)*sin(x)}
			\psStep[linecolor=teal,StepType=supremum,fillstyle=hlines,algebraic](0,9){20}{sqrt(x)*sin(x)}
			\psStep[linecolor=blue,StepType=infimum,fillstyle=vlines,algebraic](0,9){20}{sqrt(x)*sin(x)}
	\end{pspicture}}{סכומי דרבו עליונים בכחול, סכומי דרבו תחתונים בירוק, בעבור חלוקה ''אחידה``}
	
	\defi{בהינתן ההנחות מההגדרה הקודמת, נגדיר את \textit{סכום דרבו העליון} ו\textit{סכום דרבו התחתון} להיות:
	\begin{align*}
		\ol{\Sigma}(f, \Pi) = \sum_{i = 1}^{n} M_i\Delta x_i && \dard(f, \Pi) := \sum_{i = 1}^{n}m_i \Delta x_i
	\end{align*}}
	
	
	\theo{תהי $f \co [a, b] \to \R$ חסומה, ו־$\Pi$ חלוקה של $[a, b]$. אז	לכל תיוג $\{t_i\}$ של החלוקה $\Pi$, מתקיים: 
		\[ \dard(f, \Pi), S(f, \Pi, \{t_i\}) \le \daru(f, \Pi) \]
	}
	\begin{proof}
		לכל $i \in [n]$, מתקיים $m_i \le f(t_i) < M_i$. נכפיל ב־$\Delta x_i$ את הא''שים ונסכם על $i$. נקבל: 
		\[ \dard(f, \Pi) = \sumnio m_i\Delta x_i \le \sum_{i = 1}^{n}f(t_i)\Delta x_i \le \sum_{i = 1}^{n} M_i \Delta x_i = \daru(f, \Pi) \]
		כנדרש. 
	\end{proof}
	\lem{תהי $f \co [a, b] \to \R$ חסומה, וכן $\Pi$ חלוקה של $[a, b]$. אז: 
	\begin{align*}
		\daru(f, \Pi) &= \sup\ccb{S(f, \Pi, \{t_i\}) \mid \Pi \ \text{תיוג של} \ \{t_i\}} \\
		\dard(f, \Pi) &= \inf\ccb{S(f, \Pi, \{t_i\}) \mid \Pi \ \text{תיוג של} \ \{t_i\}}
	\end{align*}}
	ברשימות של שירי יש הוכחה קונקרטית עבור $\daru$, בעוד ההוכחה בעבור $\dard$ אמורה להופיע בתרגיל הבית (''אתם מעתיקים מהרשימות של שירי, וכל דבר עליון, תחתון``). לב הטיעון יושב בתוך (הכללה באינדוקציה של) תרגיל הפתיחה של ההרצאה הזו. אנחנו גם נוכיח בעבור $\daru$:
	\begin{proof}
		נגדיר $A_i = \{f(t_i) \mid t_i \in [x_{i - 1}, x_i]\}$ (כלומר $A_i = f([x_{i - 1}, x_i])$). אז: 
		\begin{multline*}
			\Sigma A := \ccb{S(f, \Pi, \{t_i\}) \mid \Pi \ \text{תיוג של} \ \{t_i\}} = \ccb{\sum_{i = 1}^{n}f(t_i)\Delta x_i \mid \forall i \co t_i \in [x_{i - 1}, x_i} = \ccb{\sumnio a_i \Delta x_i \mid a_i \in A_i} \\ 
			= (\Delta x_1)A_1 + \cdots + (\Delta x_n) A_n
		\end{multline*}
		המושגים של ''לחבר קבוצות`` ו''להכפיל קבוצה במספר`` נתונים בתרגיל הפתיחה, שם נמצא עיקר התוכן. ברשימות של שירי זה כתוב יותר מפורש. \\
		נזכר שבעצם יש להוכיח ש־$\daru(f, \Pi) = \sup \Sigma A$. מהשוויונות לעיל נקבל ש־:
		\begin{align*}
			\sup \Sigma A &= \sup((\Delta x_1)A_1 + \cdots + (\Delta x_n)A_n) \\
			&= \Delta x_1 \underbrace{\sup A_1}_{\smash{M_1}} + \cdots + \Delta x_n \underbrace{\sup A_n}_{\smash{M_n}} \\
			&= \sum_{i = 1}^{n}M_i\Delta x_i = \daru(f, \Pi)
		\end{align*}
	\end{proof}
	
	\begin{Theorem}[מונוטוניות של סכומי דרבו ביחס לעידון החלוקות]
		תהי $f \co [a, b] \to \R$ חסומה. יהיו שתי חלוקות $\Pi_1 \subset \Pi_2$ חלוקות של $[a, b]$. אז: 
		\[ \dard(f, \Pi_1) \le \dard(f, \Pi_2) \le \daru(f, \Pi_2) \le \daru(f, \Pi_1) \]
	\end{Theorem}
	אינטואיטיבית: כשמעדנים את החלוקה, סכום דרבו העליון והתחתון מתקרבים זה לזה, אך לא בהכרח משתווים. היינו רוצים שהם יתקרבו אל האינטגרל. שימו לב: חצי מהא''שים במשפט ידועים מהמשפט הקודם. 
	\begin{proof}[הוכחה עבור סכום דרבו עליון, במקרה הפרטי של $\Pi_2 = \Pi_1 \uplus \{p\}$]
		נרשום $\Pi_1 \co a = x_0 < x_1 < \cdots < x_n = b$. ישנו $k$ כך ש־$x_{k - 1} < p < x_k$. נסמן $M_i = \sup_{[x_{i - 1}, x_i]} f$, וכן $M'_k = \sup_{[x_{k - 1}, p]} f$ ו־$M''_k = \sup_{[p, x_k]}f$. אז: 
		\[ \daru(f, \Pi_2) = \sum_{i \neq k} M_i \Delta x_i  + \underbrace{\underbrace{M'_k}_{\le M_k}(p - x_{k - 1}) + \underbrace{M''_k}_{\le M_k}(x_k - p)}_{\le M_k(x_k - x_{k -1})} \le M_k\Delta x_k + \sum_{i \neq k} M_i\Delta x_i = \daru(f, \Pi_1) \]
		כנדרש. 
	\end{proof}
	הוכחת המקרה הכללי באינדוקציה על הנקודות שנוספו לחלוקה (זכרו, חלוקות הן סופיות). הוכחה עבור סכום דרבו תחתון מדואליות. 
	
	\begin{Collary}
		לכל שתי חלוקות $\Pi_1, \Pi_2$ של $[a, b]$, תמיד $\dard(f, \Pi_1) \le \daru(f, \Pi_2)$. 
	\end{Collary}
	\begin{proof}
		האיחוד $\Pi_3 := \Pi_1 \cup \Pi_2$ הוא עידון של $\Pi_1$ ושל $\Pi_2$ (עידון משותף). מהמשפט הקודם: 
		\[ \dard(f, \Pi_1) \le \dard(f, \Pi_3) \le \daru(f, \Pi_3) \le \daru(f, \Pi_2) \]
		כנדרש.
	\end{proof}
	זו מסקנה נחמדה, שאומרת שקבוצת כל סכומי הדרבו העליונים חסומה מלמטה (ע''י סכום דרבו תחתון כלשהו) ולהפך. זה אומר שקיים לה סופרמום. 
	
	\defi{תהי $f \co [a, b] \to \R$ חסומה. נגדיר: 
	\begin{itemize}
		\item אינטגרל דרבו עליון: \hfill $\displaystyle \Idaru(f) := \inf \ccb{\daru(f, \Pi) \mid [a, b] \ \text{חלוקה של} \ \Pi}$
		\item אינטגרל דרבו תחתון: \hfill $\displaystyle \Idard(f) := \sup \ccb{\dard(f, \Pi) \mid [a, b] \ \text{חלוקה של} \ \Pi}$
	\end{itemize}}

	\defi{$f$ \textit{אינטגרבילית דרבו} אם $\Idard(f) = \Idaru(f)$. אם זה קורה, $I(f) = \Idard(f) = \Idaru(f)$ נקרא \textit{אינטגרל דרבו}. }
	
	\theo{עבור $f \co [a, b]\to \R$, $f$ אינטגרבילית רימן אמ''מ $f$ חסומה וגם $f$ אינטגרבילית דרבו. אם תנאים אלו מתקיימים, אז $\int^{b}_af = I(f)$. }
	
	\rmark{אינטגרל דרבו העליון והתחתון לא מוגדרים במידה ו־$f$ איננה חסומה, ולכן זו תמיד תתווסף כהנחה למשפטים הדורשים אינטגרל דרבו עליון או תחתון. }
	
	\begin{Theorem}[קריטריון דרבו]
		$f \co [a, b] \to \R$ אינטגרבילית אמ''מ $f$ חסומה וגם לכל $\eg > 0$, כך שלכל חלוקה $\Pi$ של $[a, b]$, אם $\lg(\Pi) < \dg$ אז $\daru(f, \Pi) - \dard(f, \Pi) < \eg$. 
	\end{Theorem}
	\begin{Theorem}[קריטריון דרבו המשופר]
		$f \co [a, b] \to \R$ אינטגרבילית אמ''מ $f$ חסומה וגם לכל $\eg > 0$ קיימת חלוקה $\Pi$ של $[a, b]$ עם $\daru(f, \Pi) - \dard(f, \Pi) < \eg$. 
	\end{Theorem}
	
	\begin{proof}[הוכחת קריטריון דרבו הרגיל, החלק של ''רק אם``]
		נניח ש־$f$ אינטגרבילית רימן עם $\int^{b}_a f = I$. כבר הראינו ש־$f$ חסומה. יהי $\eg > 0$. נבחר $\hat \eg = ?$ (נבחר בהמשך את ערכו). בהתאם להגדרת אינטגרל רימן, תהי $\dg$ שמתאימה ל־$\hat \eg$. תהי $\Pi$ חלוקה עם $\lg(\Pi) < \dg$. מבחירת $\dg$, לכל תיוג $\{t_i\}$ של החלוקה, נקבל: 
		\[ I - \hat \eg < S(f, \Pi, \{t_i\}) < I + \hat \eg \]
		מטענה קודמת, ומשום שספורמום מכבד א''ש חלש, נקבל מיידית ש־: 
		\[ I - \hat \eg \le \int_{\{t_i\}} S(f, \Pi, \{t_i\}) = \dard(f, \Pi) \le \daru(f, \Pi) = \sup_{\{t_i\}} S(f, \Pi, \{t_i\}) \le I + \eg \]
		נובע: 
		\[ \daru(f, \Pi) - \dard(f, \Pi) \le 2 \hat \eg \]
		מכאן שעבור בחירת $\hat \eg = \frac{\eg}{3}$ נקבל את הנדרש. 
	\end{proof}
	הכיוון ''אם`` השני – מהרשימות של שירי. 
	
	
	
	
	
	\ndoc
\end{document}